\documentclass[aspectratio=169]{beamer}
\usetheme{Madrid}
\usecolortheme{whale}
\usepackage[utf8]{inputenc}
\usepackage[spanish]{babel}
\usepackage{amsmath,amssymb}
\usepackage{graphicx}
\usepackage{booktabs}
\usepackage{listings}
\usepackage{xcolor}

\definecolor{codegreen}{rgb}{0,0.6,0}
\definecolor{codegray}{rgb}{0.5,0.5,0.5}
\definecolor{codepurple}{rgb}{0.58,0,0.82}
\definecolor{backcolour}{rgb}{0.95,0.95,0.92}

\lstdefinestyle{pythonstyle}{
    backgroundcolor=\color{backcolour},
    commentstyle=\color{codegreen},
    keywordstyle=\color{blue},
    numberstyle=\tiny\color{codegray},
    stringstyle=\color{codepurple},
    basicstyle=\ttfamily\footnotesize,
    breakatwhitespace=false,
    breaklines=true,
    captionpos=b,
    keepspaces=true,
    showspaces=false,
    showstringspaces=false,
    showtabs=false,
    tabsize=2
}

\lstset{style=pythonstyle}

\title{Correlación y ANCOVA}
\subtitle{Relación entre variables y análisis de covarianza}
\author{Métodos de Análisis de Datos 1}
\institute{Universidad Nacional del Sur}
\date{Unidad 3 - Clase 19}

\begin{document}

\frame{\titlepage}

\begin{frame}{Contenidos}
\tableofcontents
\end{frame}

\section{Correlación}

\begin{frame}{¿Qué es la correlación?}
\begin{block}{Definición}
La correlación mide la \textbf{fuerza y dirección} de la relación lineal entre dos variables.
\end{block}

\vspace{0.5cm}
\textbf{Diferencias con regresión:}
\begin{itemize}
    \item Correlación: simétrica ($r_{XY} = r_{YX}$)
    \item Regresión: asimétrica (hay variable respuesta y predictora)
    \item Correlación: solo mide asociación
    \item Regresión: permite predicción
\end{itemize}

\vspace{0.3cm}
\textbf{Relación:} En regresión simple, $R^2 = r_{XY}^2$
\end{frame}

\begin{frame}{Correlación de Pearson}
\begin{block}{Coeficiente de correlación de Pearson}
\[
r_{XY} = \frac{\sum_{i=1}^{n}(X_i - \bar{X})(Y_i - \bar{Y})}{\sqrt{\sum_{i=1}^{n}(X_i - \bar{X})^2} \sqrt{\sum_{i=1}^{n}(Y_i - \bar{Y})^2}}
\]
\end{block}

\vspace{0.3cm}
\textbf{Propiedades:}
\begin{itemize}
    \item $-1 \leq r \leq 1$
    \item $r = 1$: correlación positiva perfecta
    \item $r = -1$: correlación negativa perfecta
    \item $r = 0$: no hay correlación \textbf{lineal}
    \item Adimensional (no depende de unidades)
\end{itemize}
\end{frame}

\begin{frame}{Interpretación}
\textbf{Reglas prácticas:}
\begin{itemize}
    \item $|r| < 0.3$: correlación débil
    \item $0.3 \leq |r| < 0.7$: correlación moderada
    \item $|r| \geq 0.7$: correlación fuerte
\end{itemize}

\vspace{0.5cm}
\textbf{Importante:} El umbral depende del contexto
\begin{itemize}
    \item En ciencias sociales: $r = 0.3$ puede ser notable
    \item En física: $r < 0.9$ puede ser insuficiente
\end{itemize}
\end{frame}

\begin{frame}{Test de significancia}
¿Es $r$ significativamente distinto de 0?

\begin{block}{Hipótesis}
\[
H_0: \rho = 0 \quad \text{vs} \quad H_1: \rho \neq 0
\]
\end{block}

\textbf{Estadístico:}
\[
t = \frac{r\sqrt{n-2}}{\sqrt{1-r^2}} \sim t_{n-2}
\]

\vspace{0.5cm}
\textbf{Supuestos:}
\begin{itemize}
    \item $(X, Y)$ siguen distribución normal bivariada
    \item Relación lineal
\end{itemize}
\end{frame}

\begin{frame}{Correlación de Spearman}
\textbf{Uso:} Cuando los datos no son normales o la relación es monótona pero no lineal

\vspace{0.3cm}
\begin{block}{Coeficiente de Spearman}
\[
r_S = 1 - \frac{6\sum_{i=1}^{n} d_i^2}{n(n^2-1)}
\]
\end{block}

donde $d_i$ es la diferencia entre los rangos de $X_i$ e $Y_i$.

\vspace{0.5cm}
\textbf{Ventajas:}
\begin{itemize}
    \item Robusto a outliers
    \item No requiere normalidad
    \item Detecta relaciones monótonas no lineales
\end{itemize}
\end{frame}

\begin{frame}{Correlación $\neq$ Causalidad}
\begin{alertblock}{Regla de oro}
\textbf{Correlación NO implica causalidad}
\end{alertblock}

\vspace{0.5cm}
\textbf{Posibles explicaciones para $r_{XY} \neq 0$:}
\begin{enumerate}
    \item $X$ causa $Y$
    \item $Y$ causa $X$
    \item Existe variable $Z$ que causa tanto $X$ como $Y$ (confusor)
    \item Correlación espuria (coincidencia)
\end{enumerate}

\vspace{0.5cm}
\textbf{Ejemplo clásico:}
\begin{itemize}
    \item Alta correlación entre ventas de helados y ahogamientos
    \item No es que los helados causen ahogamientos
    \item Ambos aumentan en verano (variable confusora)
\end{itemize}
\end{frame}

\section{ANCOVA}

\begin{frame}{¿Qué es ANCOVA?}
\textbf{ANCOVA = Análisis de Covarianza}

\vspace{0.3cm}
\begin{block}{Definición}
Combina ANOVA con regresión lineal. Compara grupos mientras se ajusta por variables continuas (covariables).
\end{block}

\vspace{0.5cm}
\textbf{Modelo:}
\[
Y_{ij} = \mu + \alpha_i + \beta X_{ij} + \epsilon_{ij}
\]

donde:
\begin{itemize}
    \item $Y_{ij}$: respuesta de observación $j$ en grupo $i$
    \item $\alpha_i$: efecto del grupo $i$
    \item $X_{ij}$: covariable continua
    \item $\beta$: pendiente (común a todos los grupos)
\end{itemize}
\end{frame}

\begin{frame}{¿Cuándo usar ANCOVA?}
\textbf{Situaciones comunes:}
\begin{enumerate}
    \item Comparar tratamientos ajustando por medida basal
    \begin{itemize}
        \item Ejemplo: Comparar dietas ajustando por peso inicial
    \end{itemize}
    
    \item Reducir variabilidad no explicada
    \begin{itemize}
        \item Ejemplo: Comparar métodos de enseñanza ajustando por nivel previo
    \end{itemize}
    
    \item Controlar por variables confusoras
    \begin{itemize}
        \item Ejemplo: Comparar salarios por género ajustando por experiencia
    \end{itemize}
\end{enumerate}

\vspace{0.5cm}
\textbf{Ventaja principal:} Mayor potencia estadística (menor error)
\end{frame}

\begin{frame}{Supuestos de ANCOVA}
Además de los supuestos de regresión lineal:

\begin{enumerate}
    \item \textbf{Independencia de covariable y tratamiento}
    \begin{itemize}
        \item La covariable no debe estar afectada por el grupo
        \item Ejemplo: peso inicial debe ser pre-tratamiento
    \end{itemize}
    
    \item \textbf{Pendientes paralelas (homogeneidad de regresión)}
    \begin{itemize}
        \item La relación $X$-$Y$ es la misma en todos los grupos
        \item Mismo $\beta$ para todos los grupos
    \end{itemize}
    
    \item \textbf{Linealidad}
    \begin{itemize}
        \item Relación lineal entre covariable y respuesta
    \end{itemize}
\end{enumerate}
\end{frame}

\begin{frame}{Test de pendientes paralelas}
\textbf{¿Cómo verificar?} Ajustar modelo con interacción:

\[
Y_{ij} = \mu + \alpha_i + \beta X_{ij} + (\alpha\beta)_i X_{ij} + \epsilon_{ij}
\]

\vspace{0.3cm}
\begin{block}{Hipótesis}
\[
H_0: (\alpha\beta)_1 = (\alpha\beta)_2 = \cdots = 0 \quad \text{(pendientes iguales)}
\]
\end{block}

\vspace{0.3cm}
\textbf{Decisión:}
\begin{itemize}
    \item Si NO rechazamos $H_0$ $\Rightarrow$ OK, usar ANCOVA
    \item Si rechazamos $H_0$ $\Rightarrow$ pendientes diferentes, ANCOVA no apropiado
\end{itemize}
\end{frame}

\begin{frame}{Interpretación de resultados}
Después de ajustar por la covariable:

\vspace{0.3cm}
\textbf{Efecto de la covariable:}
\begin{itemize}
    \item $\hat{\beta}$: cambio en $Y$ por unidad de $X$, común a todos los grupos
\end{itemize}

\vspace{0.3cm}
\textbf{Efecto del grupo:}
\begin{itemize}
    \item $\hat{\alpha}_i$: diferencias entre grupos \textbf{ajustadas} por $X$
    \item Son las diferencias que quedan después de controlar por $X$
\end{itemize}

\vspace{0.5cm}
\textbf{Ventaja:} Si la covariable explica mucha variabilidad, el error residual disminuye $\Rightarrow$ mayor potencia para detectar diferencias entre grupos
\end{frame}

\section{Implementación}

\begin{frame}[fragile]{Python: Correlación}
\begin{lstlisting}[language=Python]
import scipy.stats as stats
import numpy as np

# Correlacion de Pearson
r_pearson, p_pearson = stats.pearsonr(X, Y)
print(f"Pearson: r = {r_pearson:.3f}, p = {p_pearson:.4f}")

# Correlacion de Spearman
r_spearman, p_spearman = stats.spearmanr(X, Y)
print(f"Spearman: r_s = {r_spearman:.3f}, p = {p_spearman:.4f}")

# Matriz de correlacion
import pandas as pd
df = pd.DataFrame({'X1': X1, 'X2': X2, 'X3': X3, 'Y': Y})
corr_matrix = df.corr()
print(corr_matrix)
\end{lstlisting}
\end{frame}

\begin{frame}[fragile]{Python: ANCOVA}
\begin{lstlisting}[language=Python]
from statsmodels.formula.api import ols
from statsmodels.stats.anova import anova_lm

# Modelo ANCOVA
modelo_ancova = ols('peso_final ~ C(dieta) + peso_inicial',
                    data=datos).fit()

# Tabla ANOVA tipo II (ajustada)
print(anova_lm(modelo_ancova, typ=2))

# Test de pendientes paralelas (interaccion)
modelo_interaccion = ols('peso_final ~ C(dieta) * peso_inicial',
                         data=datos).fit()
print(anova_lm(modelo_interaccion, typ=2))

# Si interaccion NO significativa => OK usar ANCOVA
\end{lstlisting}
\end{frame}

\begin{frame}[fragile]{R: Correlación}
\begin{lstlisting}[language=R]
# Correlacion de Pearson
cor.test(X, Y, method = "pearson")

# Correlacion de Spearman
cor.test(X, Y, method = "spearman")

# Matriz de correlacion
datos <- data.frame(X1, X2, X3, Y)
cor(datos)

# Visualizacion
library(corrplot)
corrplot(cor(datos), method = "circle")
\end{lstlisting}
\end{frame}

\begin{frame}[fragile]{R: ANCOVA}
\begin{lstlisting}[language=R]
# Modelo ANCOVA
modelo_ancova <- lm(peso_final ~ dieta + peso_inicial,
                    data = datos)

# Tabla ANOVA
anova(modelo_ancova)

# O con car para tipo II
library(car)
Anova(modelo_ancova, type = "II")

# Test de pendientes paralelas
modelo_interaccion <- lm(peso_final ~ dieta * peso_inicial,
                         data = datos)
anova(modelo_interaccion)

# Comparar modelos
anova(modelo_ancova, modelo_interaccion)
\end{lstlisting}
\end{frame}

\begin{frame}[fragile]{Ejemplo completo: ANCOVA}
\textbf{Datos:} Efecto de 3 dietas sobre pérdida de peso

\begin{lstlisting}[language=Python]
# Cargar datos
dieta A: peso_inicial=[80,78,85], peso_final=[72,68,75]
dieta B: peso_inicial=[82,75,88], peso_final=[70,65,78]
dieta C: peso_inicial=[81,79,84,76], peso_final=[71,69,73,67]

# ANCOVA
modelo = ols('peso_final ~ C(dieta) + peso_inicial',
             data=datos).fit()
print(anova_lm(modelo, typ=2))

# Resultados:
#                    sum_sq    df         F    PR(>F)
# C(dieta)            120.5   2.0      8.50   0.0089
# peso_inicial        890.2   1.0     62.50   0.0001
# Residual             99.9   7.0       NaN      NaN

# Conclusion: Hay diferencias entre dietas (p=0.009)
# ajustando por peso inicial
\end{lstlisting}
\end{frame}

\section{Resumen}

\begin{frame}{Resumen de la Unidad 3}
\textbf{Clase 14-15: Regresión lineal simple}
\begin{itemize}
    \item Modelo, estimación (MCO), $R^2$
    \item Inferencia: tests, IC, predicción
\end{itemize}

\textbf{Clase 16: Diagnóstico}
\begin{itemize}
    \item Gráficos de residuos
    \item Tests formales
    \item Medidas de influencia
\end{itemize}

\textbf{Clase 17: Regresión múltiple}
\begin{itemize}
    \item Modelo con múltiples predictores
    \item Interpretación ajustada
    \item Multicolinealidad (VIF)
\end{itemize}
\end{frame}

\begin{frame}{Resumen de la Unidad 3 (cont.)}
\textbf{Clase 18: Selección y polinomios}
\begin{itemize}
    \item Criterios: AIC, BIC
    \item Métodos: forward, backward, stepwise
    \item Regresión polinomial
\end{itemize}

\textbf{Clase 19: Correlación y ANCOVA}
\begin{itemize}
    \item Correlación de Pearson y Spearman
    \item Correlación $\neq$ causalidad
    \item ANCOVA: comparar grupos ajustando por covariable
\end{itemize}

\vspace{0.5cm}
\textbf{Próxima unidad:} ANOVA y diseño de experimentos
\end{frame}

\end{document}
